\begin{figure}[ht]
\centering
\begin{tikzpicture}[>=Stealth, font=\small, node distance=1.35cm]
  \node[draw, rounded corners, fill=blue!8, minimum width=4.2cm, minimum height=0.95cm] (gq)
    {profinite 群与 Frobenius};
  \node[draw, rounded corners, fill=green!10, below left=1.7cm and 2.4cm of gq, minimum width=4.2cm, minimum height=0.95cm] (cyc)
    {分圆特征与一维表示};
  \node[draw, rounded corners, fill=green!10, below right=1.7cm and 2.4cm of gq, minimum width=4.2cm, minimum height=0.95cm] (tate)
    {Tate 模与 Weil 配对};
  \node[draw, rounded corners, fill=orange!12, below=2.8cm of gq, minimum width=4.8cm, minimum height=0.95cm] (mod)
    {模形式表示与水平下降实例};

  \draw[->, very thick] (gq) -- (cyc);
  \draw[->, very thick] (gq) -- (tate);
  \draw[->, very thick] (cyc) -- (mod);
  \draw[->, very thick] (tate) -- (mod);
\end{tikzpicture}
\caption{习题主线：先熟悉 Frobenius、分圆特征与 Tate 模，再把这些概念推到模形式表示与水平下降的数值例子上。}
\label{fig:ch09-exercises-roadmap}
\end{figure}
